\documentclass[12pt, ukrainian]{article}
\usepackage[a4paper]{geometry}
\setlength\parindent{0mm}
\geometry{top=1.1cm,bottom=1.1cm,left=1.1cm,right=1.1cm}
\pagestyle{empty}
\usepackage[T2A]{fontenc}
\usepackage[utf8]{inputenc}
\usepackage[ukrainian]{babel}
\usepackage{graphicx}
\usepackage{caption}
\usepackage{verbatim, listings, clrscode3e, fancyvrb}
\DeclareCaptionFormat{nocaption}{}
\captionsetup{format=nocaption,aboveskip=0pt,belowskip=0pt}
\usepackage[Export]{adjustbox}
\adjustboxset{max size={0.9\linewidth}{0.9\paperheight}}
\def\code#1{\Verb!#1!}
\begin{document}
{{16.05.2024 Контрольна робота, варіант  \No \textbf { 32 }} група К\_\_ ПІБ {\parbox {2.5 cm} \dotfill}\par}
{
%\begin {large}
\begin{minipage}[t]{9.7cm}
У завданнях 1--3 вказати, що буде виведено за виконання. Повідомлення інтерпретатора про помилку позначати ERROR


               
{\begin{center}Завдання {1}\end{center}}
\vspace{-4mm}

\begin{verbatim}
d = ('a','b','c','d','e',0,1,2,3,4)
print(d[14])

\end{verbatim}            

\vspace{15mm}
               
{\begin{center}Завдання {2}\end{center}}
\vspace{-4mm}

\begin{verbatim}
try:
    t = {25:6, 44:3, 55:7, 55:8}
    t[55] = 3
    for x in t.keys():
        print(x, sep=' ', end=' ')
except:
    print('ERR')

\end{verbatim}            


\end{minipage}
\vline \hspace{8mm}
\begin{minipage}[t]{8.0cm}                
               
{\begin{center}Завдання {3}\end{center}}
\vspace{-4mm}

\begin{verbatim}
d = 0
n = 1

class C:
    n = 2
    
    def __init__(self):
        global n
        self.n = 3
        n = 4
        C.n = 5

obj = C()
try:
    print(d, end=' ')
    print(n, end=' ')
    print(C.n, end=' ')
    print(obj.n)
except:
    print('ERR')

\end{verbatim}            

\end{minipage}


               
{\begin{center}Завдання {4}\end{center}}
\vspace{-4mm}

Записати ВИРАЗ, значенням якого є список, що складається з цілих чисел від 20 до 2005 (включно), що не діляться на жодне з чисел 2, 3.\par

\vspace{10mm}
               
{\begin{center}Завдання {5}\end{center}}
\vspace{-4mm}
Змінні \kw{next} та \kw{prev} вузла списку позначають наступний та попередній за ним вузол відповідно. Починаючи з голови, у список записано послідовні цілі числа від~$-64$ до~$40$. Нехай змінна \kw{p} позначає вузол, що зберігає значення~$2$.
Яке число зберігається у вузлі \kw{p.prev.prev.next.prev.prev} ?\par

               
{\begin{center}Завдання {6}\end{center}}
\vspace{-4mm}
Значенням змінної \kw{a} є цілочисловий масив типу $numpy.ndarray$ розмірів $7{\times}4$. На скільки збільшиться сума елементів масиву \kw{a} після виконання
\begin{verbatim}
a[2] += 3
\end{verbatim}


Якщо код некоректний, то в якості відповіді зазначити ERROR.\par

               
{\begin{center}Завдання {10}\end{center}}
\vspace{-4mm}

Написати функцію $f$, яка для файлу, заданого шляхом, розв’язує задачу:\\ Текстовий файл містить послідовність чисел $a_1$, $a_2$, $\ldots$, $a_t$, записану у форматі одне число на рядок. Знайти середнє арифметичне чисел $a_1$, $a_4$, $a_7$, $\ldots$ (починаючи з першого через два).

В усіх випадках розміри файлів такі, що їх вміст повністю завантажений у пам'ять бути не може. Якість коду також оцінюється.



\vspace{10mm}\par

\newpage
\begin{minipage}[t]{9.7cm}
               
{\begin{center}Завдання {7}\end{center}}
\vspace{-4mm}

\begin{center}
	\adjustimage{max size={0.8\linewidth}{0.8\paperheight}}
	{../15BinTree/trees_exp/104.png}
\end{center}
Указати послідовність міток вершин дерева, зображеного на рис., що утвориться в результаті обходу дерева в прямому порядку. Мітки записувати через пробіл.\par Алгоритм починає роботу з кореня (на рис. це верхня вершина).\par

\end{minipage}
\vline \hspace{8mm}
\begin{minipage}[t]{8.0cm}
               
{\begin{center}Завдання {8}\end{center}}
\vspace{-4mm}
Для графа на рисунку з вершини 3 запущено алгоритм пошуку в глибину, що будує кістякове дерево. Списки суміжності вершин переглядаються алгоритмом за зростанням міток вершин.\par
\begin{center}
	\adjustimage{max size={0.9\linewidth}{0.9\paperheight}}
	{../16Graphs/Traversals/93.png}
\end{center}
\par Які з наведених ребер входять до складу отриманого кістякового дерева?\par
1) (1, 3)\par
2) (1, 5)\par
3) (0, 2)\par
4) (3, 5)\par
5) (3, 4)\par
6) (2, 4)\par
{\vskip 15pt} \par

\end{minipage}

               
{\begin{center}Завдання {9}\end{center}}
\vspace{-4mm}
Файл \code{'1.txt'} містить журнал з рядками, в яких через двокрапку записано поряд\-ко\-вий номер запису (ціле), ім'я користувача (складається з латиниці), час події (фор\-мат як у прикладі), тип події (ERROR, WARNING, INFO, STATUS), код події (ціле).

Білі символи навколо роздільників не допускаються.

Приклад рядка журналу



\begin{verbatim}
12:username:2022-05-05 13-26:ERROR:404
\end{verbatim}


Нижче наведено код, який шукає усі попередження (подія WARNING), що зустрі\-ча\-ються в журналі і відбулися під час обідньої перерви деканату (з 13:00 по 13:59 включно), та зберігає їх коди в змінній \kw{codes}.



\begin{verbatim}
codes = set()
regex = re.compile(r'XXX')
with open('1.txt', encoding='utf8') as f:
    for s in f:
        res = re.fullmatch(regex, s)
        if not res: continue
        s = ''
        codes.add(YYY)
\end{verbatim}





Який рядок треба записати на місце \kw{XXX} ?
Який рядок треба записати на місце \kw{YYY} ?\par

%\end {large}
}
%end
\end{document}

